\chapter{Hyperactivity}

\section*{Definition}
Hyperactivity refers to excessive, often purposeless motor activity such as fidgeting, restlessness, running or climbing in inappropriate settings, and an inability to remain seated. In children, it is a core symptom of attention-deficit/hyperactivity disorder (ADHD) when present for at least 6 months across multiple settings.

\section*{What Happens in Your Body}
Hyperactivity involves underactivation of prefrontal-striatal-cerebellar circuits that regulate motor inhibition and attention. Hypofunctioning dopaminergic and noradrenergic signaling reduces the ability to suppress motor output, while altered dopamine transporter (DAT) density in the striatum impairs reward anticipation and response control.

\section*{Causes}
\begin{itemize}
  \item \textbf{Psychiatric:} Attention-deficit/hyperactivity disorder (ADHD), bipolar disorder (manic episode), anxiety disorders, autism spectrum disorder
  \item \textbf{Neurological:} Tourette syndrome, frontal lobe lesions, neuroacanthocytosis, Kleine-Levin syndrome
  \item \textbf{Endocrine/metabolic:} Hyperthyroidism, hypoglycemia, pheochromocytoma
  \item \textbf{Substance-induced:} Caffeine, nicotine, stimulants, corticosteroids, bronchodilators, antiepileptics (levetiracetam)
  \item \textbf{Psychosocial:} Inadequate sleep, chaotic environment, child abuse or neglect
\end{itemize}

\section*{When to See a Doctor}
See your doctor if:
\begin{itemize}
  \item Hyperactivity negatively impacts school, work, or social functioning
  \item Onset is sudden or associated with tics, sleep disturbance, or physical symptoms
  \item There is concern for bipolar disorder (euphoria, decreased need for sleep) or hyperthyroidism
\end{itemize}

\section*{Related Symptoms}
\begin{itemize}
  \item Impulsivity, inattention, restlessness, fidgeting, tics, pressured speech, distractibility, sleep disturbance
\end{itemize}
\textbf{Important:} Normal age-appropriate activity levels in children must be distinguished from pathological hyperactivity by assessing pervasiveness and impairment.

\section*{Self-Care / Home Management}
\begin{itemize}
  \item Ensure regular physical activity to channel excess energy
  \item Maintain consistent sleep and mealtime routines
  \item Minimize sugar and avoid caffeine
  \item Use organizational tools (lists, timers, visual schedules) to structure tasks
\end{itemize}

\section*{How Your Doctor Will Evaluate You}
\begin{itemize}
  \item Clinical interview and DSM-5 criteria for ADHD; behavior rating scales (Conners, Vanderbilt, SNAP-IV)
  \item Collateral information from teachers, parents, or partners
  \item Thyroid function tests, blood glucose, toxicology screen if environmental factors suspected
  \item Neuropsychological testing if diagnosis is unclear or comorbid conditions suspected
\end{itemize}

\section*{Treatment Options}
\begin{itemize}
  \item First-line for ADHD: stimulants (methylphenidate, amphetamine salts) or non-stimulants (atomoxetine, guanfacine, clonidine)
  \item Behavioral therapy: parent training in behavioral management, organizational skills training, cognitive-behavioral therapy
  \item Classroom or workplace accommodations (preferred seating, breaks, reduced distractions)
  \item Treat underlying causes: antithyroid medication for hyperthyroidism, mood stabilizers for bipolar disorder
\end{itemize}

\section*{References}
\begin{thebibliography}{99}
\bibitem{hyperactivity:ref1} American Psychiatric Association. ``Diagnostic and Statistical Manual of Mental Disorders.'' 5th ed, Text Revision. APA, 2022.
\bibitem{hyperactivity:ref2} Faraone SV, Asherson P, Banaschewski T, et al. ``Attention-Deficit/Hyperactivity Disorder.'' Nat Rev Dis Primers. 2015;1:15020.
\bibitem{hyperactivity:ref3} Wolraich ML, Hagan JF Jr, Allan C, et al. ``Clinical Practice Guideline for the Diagnosis, Evaluation, and Treatment of ADHD in Children and Adolescents.'' Pediatrics. 2019;144(4):e20192528.
\bibitem{hyperactivity:ref4} Pliszka S; AACAP Work Group on Quality Issues. ``Practice Parameter for the Assessment and Treatment of Children and Adolescents with ADHD.'' J Am Acad Child Adolesc Psychiatry. 2007;46(7):894--921.
\bibitem{hyperactivity:ref5} Cortese S, Adamo N, Del Giovane C, et al. ``Comparative Efficacy and Tolerability of Medications for ADHD in Children and Adolescents.'' Lancet Psychiatry. 2018;5(9):727--738.
\bibitem{hyperactivity:ref6} Kliegman RM, St Geme JW III, Blum NJ, et al. ``Nelson Textbook of Pediatrics.'' 22nd ed. Elsevier; 2024.
\end{thebibliography}
